\section{AHSME 1951}
        \begin{enumerate}
        	\item (\textit{AHSME 1951}) The percent that $M$ is greater than $N$ is:


 $\textbf{(A) \ } \frac {100(M - N)}{M} \qquad \textbf{(B) \ } \frac {100(M - N)}{N} \qquad \textbf{(C) \ } \frac {M - N}{N} \qquad \textbf{(D) \ } \frac {M - N}{M} \qquad \textbf{(E) \ } \frac {100(M + N)}{N}$

 \textbf{Answer: }B



\textbf{Solution}

$M-N$ is the amount by which $M$ is greater than $N$. We divide this by $N$ to get the percent by which $N$ is increased in the form of a decimal, and then multiply by $100$ to make it a percentage. Therefore, the answer is $\boxed{\mathrm{(B)}\ \dfrac{100(M-N)}{N}}$.

 


 


	\item (\textit{AHSME 1951}) A rectangular field is half as wide as it is long and is completely enclosed by $x$ yards of fencing. The area in terms of $x$ is:


 $(\textbf{A})\ \frac{x^2}2 \qquad (\textbf{B})\ 2x^2 \qquad (\textbf{C})\ \frac{2x^2}9 \qquad (\textbf{D})\ \frac{x^2}{18} \qquad (\textbf{E})\ \frac{x^2}{72}$

 \textbf{Answer: }D



\textbf{Solution}

Let $w$ be the width. Then $l = 2w$, and the perimeter is $x = 2(2w)+2w = 6w \implies w = \frac{x}6$. The area is $wl = w(2w) = 2w^2 = 2\left(\frac{x^2}{36}\right) = \frac{x^2}{18}$, so the answer is $\boxed{\textbf{(D)}}$



	\item (\textit{AHSME 1951}) If the length of a diagonal of a square is $a + b$, then the area of the square is:


 $\textbf{(A)}\ (a+b)^{2}\qquad\textbf{(B)}\ \frac{1}{2}(a+b)^{2}\qquad\textbf{(C)}\ a^{2}+b^{2}$
$\textbf{(D)}\ \frac{1}{2}(a^{2}+b^{2})\qquad\textbf{(E)}\ \text{none of these}$

 \textbf{Answer: }B



\textbf{Solution 1}

Let $s$ be the side length of the square. By the \href{https://artofproblemsolving.com/wiki/index.php?title=Pythagorean\_Theorem}{Pythagorean Theorem}, $s^2 + s^2 = 2s^2 = (a+b)^2$, so the area of the square is $s^2 = \boxed{\textbf{(B)}\  \frac{1}{2}(a+b)^2}$.

 



\textbf{Solution 2}

Recall that the area of a rhombus is equal to half the product of its diagonals, so the area of this square is $\boxed{\textbf{(B)}\  \frac{1}{2}(a+b)^2}$.

 


	\item (\textit{AHSME 1951}) A barn with a flat roof is rectangular in shape, $10$ yd. wide, $13$ yd. long and $5$ yd. high.  It is to be painted inside and outside, and on the ceiling, but not on the roof or floor. The total number of sq. yd. to be painted is:


 $\textbf{(A) \ } 360 \qquad \textbf{(B) \ } 460 \qquad \textbf{(C) \ } 490 \qquad \textbf{(D) \ } 590 \qquad \textbf{(E) \ } 720$

 \textbf{Answer: }D



\textbf{Solution}

The walls are $13*5=65$ and $10*5=50$ in area, and the ceiling has an area of $10*13=130$.

 $((65+50)2)2+130=590 \Rightarrow \boxed{\mathrm{(D)}}$



	\item (\textit{AHSME 1951}) Mr. A owns a home worth $\$10,000$.  He sells it to Mr. B at a $10 \%$ profit based on the worth of the house. Mr. B sells the house back to Mr. A at a $10 \%$ loss.  Then:


 $\textbf{(A)}\ \text{A comes out even} \qquad\textbf{(B)}\ \text{A makes \$1100 on the deal} \qquad\textbf{(C)}\ \text{A makes \$1000 on the deal}$
$\textbf{(D)}\ \text{A loses \$900 on the deal} \qquad\textbf{(E)}\ \text{A loses \$1000 on the deal}$

 \textbf{Answer: }B



\textbf{Solution}

Mr. A sells his home to Mr. B for $(1 + 0.10) \cdot \$10000 = \$11000$. Then, Mr. B sells it back to Mr. A for $(1 - 0.10) \cdot \$11000 = \$9900$. Therefore, Mr. A makes $\$11000-\$9900=\$1100$ on the deal. $\boxed{\textbf{(B)}}$



	\item (\textit{AHSME 1951}) The bottom, side, and front areas of a rectangular box are known. The product of these areas is equal to:


 $\textbf{(A)}\ \text{the volume of the box} \qquad\textbf{(B)}\ \text{the square root of the volume} \qquad\textbf{(C)}\ \text{twice the volume}$
$\textbf{(D)}\ \text{the square of the volume} \qquad\textbf{(E)}\ \text{the cube of the volume}$

 \textbf{Answer: }D



\textbf{Solution}

Let $a$, $b$, and $c$ be the dimensions of the box. Then the areas of the aforementioned faces $ab$, $bc$, and $ac$, while the volume of the box is $abc$. The product of these three areas is $ab \cdot bc \cdot ac = a^2b^2c^2 = (abc)^2$, which is $\boxed{\textbf{(D)}\ \text{the square of the volume}}$ of the box.

 


	\item (\textit{AHSME 1951}) An error of $.02"$ is made in the measurement of a line $10"$ long, while an error of only $.2"$ is made in a measurement of a line $100"$ long. In comparison with the relative error of the first measurement, the relative error of the second measurement is:


 $\textbf{(A) \ } \text{greater by }.18 \qquad\textbf{(B) \ } \text{the same} \qquad \textbf{(C) \ } \text{less} \qquad\textbf{(D) \ } 10\text{ times as great} \qquad\textbf{(E) \ } \text{correctly described by both}$

 \textbf{Answer: }B



\textbf{Solution}

The relative error of the first measurement is $\frac{0.02}{10} = 0.002$, while the relative error of the second measurement is $\frac{0.2}{100} = 0.002$. Therefore, these relative errors are $\boxed{\textbf{(B)}\ \text{the same}}$.

 


	\item (\textit{AHSME 1951}) The price of an article is cut $10 \%.$ To restore it to its former value, the new price must be increased by:


 $\textbf{(A) \ } 10 \% \qquad\textbf{(B) \ } 9 \% \qquad \textbf{(C) \ } 11\frac{1}{9} \% \qquad\textbf{(D) \ } 11 \% \qquad\textbf{(E) \ } \text{none of these answers}$

 \textbf{Answer: }C



\textbf{Solution}

\href{https://artofproblemsolving.com/wiki/index.php?title=Without\_loss\_of\_generality}{Without loss of generality}, let the price of the article be 100. Thus, the new price is $.9\cdot  100= 90$. Then, to restore it to the original price, we solve $90x=100$. We find $x=1\dfrac{1}{9}$, thus, the percent increase is $1\dfrac{1}{9}-1=\dfrac{1}{9}=\boxed{\textbf{(C) \ } 11\frac{1}{9}\%}$.

 


	\item (\textit{AHSME 1951}) An equilateral triangle is drawn with a side length of $a.$ A new equilateral triangle is formed by joining the midpoints of the sides of the first one. then a third equilateral triangle is formed by joining the midpoints of the sides of the second; and so on forever. the limit of the sum of the perimeters of all the triangles thus drawn is:


 $\textbf{(A) \ } \text{Infinite} \qquad\textbf{(B) \ } 5\frac{1}{4}a \qquad \textbf{(C) \ } 2a \qquad\textbf{(D) \ } 6a \qquad\textbf{(E) \ } 4\frac{1}{2}a$

 \textbf{Answer: }D



\textbf{Solution}

The perimeter of the first triangle is $3a$. By the midsegment theorem, the side length of any triangle after the first must be half the side length of the previous triangle, so the perimeter of any triangle after the first must be half the perimeter of the previous triangle. 

 Therefore, the sum of the perimeters of all the triangles is the infinite geometric series $3a+\frac{3a}{2}+\frac{3a}{4}+\cdots$.  This series' first term is $3a$ while its common ratio is $\frac{1}{2}$, so its sum is $\frac{3a}{1-\frac{1}{2}}=\boxed{\textbf{(D)}\ 6a}$.

 


	\item (\textit{AHSME 1951}) Of the following statements, the one that is incorrect is:


 $\textbf{(A)}\ \text{Doubling the base of a given rectangle doubles the area.}\qquad\textbf{(B)}\ \text{Doubling the altitude of a triangle doubles the area.}$
$\textbf{(C)}\ \text{Doubling the radius of a given circle doubles the area.}$
$\textbf{(D)}\ \text{Doubling the divisor of a fraction and dividing its numerator by 2 changes the quotient.}$
$\textbf{(E)}\ \text{Doubling a given quantity may make it less than it originally was.}$

 \textbf{Answer: }C



\textbf{Solution}

The well-known area formula for a circle is $A = \pi r^2$, so doubling the radius will result it quadrupling the area (since $A' = \pi (2r)^2 = 4 \pi r^2 = 4A$). Statement $\boxed{\textbf{(C)}}$ is therefore incorrect, and is the correct answer choice.

 
Statements $\textbf{(A)}$ and $\textbf{(B)}$ are evidently correct, since in triangles the area is directly proportional to both the base and the height ($A = \frac{1}{2} bh$). Statement $\textbf{(D)}$ is also correct: let $q = \frac{a}{b}$. Then $q' = \frac{a \div 2}{2b} = \frac{a}{4b} = \frac{q}{4}$, which is a change from $q$. Finally, statement $\textbf{(E)}$ is true: $2x < x$ if $x < 0$, as doubling a negative number makes it even more negative (and therefore less than it originally was).

 


	\item (\textit{AHSME 1951}) The limit of the sum of an infinite number of terms in a geometric progression is $\frac {a}{1- r}$ where $a$ denotes the first term and $-1 < r < 1$ denotes the common ratio. The limit of the sum of their squares is:


 $\textbf{(A)}\ \frac {a^2}{(1 -r)^2} \qquad\textbf{(B)}\ \frac {a^2}{1 + r^2} \qquad\textbf{(C)}\ \frac {a^2}{1 - r^2} \qquad\textbf{(D)}\ \frac {4a^2}{1+ r^2} \qquad\textbf{(E)}\ \text{none of these}$

 \textbf{Answer: }C



\textbf{Solution}

The original series is $a + ar + ar^2 + ar^3 + \cdots$, so the new series is $a^2 + a^2r^2 +a^2r^4 +a^2r^6+\cdots$.  The new series is an infinite geometric series with first term $a^2$ and common ratio $r^2$, so its sum is $\boxed{\textbf{(C)}\ \frac {a^2}{1 - r^2}}$.

 


	\item (\textit{AHSME 1951}) At $2: 15$ o'clock, the hour and minute hands of a clock form an angle of:


 $\textbf{(A)}\ 30^{\circ} \qquad\textbf{(B)}\ 5^{\circ} \qquad\textbf{(C)}\ 22\frac {1}{2}^{\circ} \qquad\textbf{(D)}\ 7\frac {1}{2} ^{\circ} \qquad\textbf{(E)}\ 28^{\circ}$

 \textbf{Answer: }C



\textbf{Solution 1}

At $2:15$, the minute hand points directly at $3$, while the hour hand points $\frac{1}{4}$ of the way from $2$ to $3$.  The angle between any two numbers on the clock is $\frac{360^{\circ}}{12} = 30^{\circ}$, so the angle between the hands at $2:15$ is $\frac{3}{4} \cdot 30^{\circ} = \boxed{\textbf{(C)}\ 22\frac{1}{2}^{\circ}}$.

 



\textbf{Solution 2}

Using the formula $\frac{|60h-11m|}{2}$, where $h$ is the hour time and $m$ is the minute time, we get that the angle is $\frac{|60(2)-11(15)|}{2}=\frac{|120-165|}{2}=\frac{45}{2}=\boxed{\textbf{(C)}\ 22\frac {1}{2}^{\circ} }$



	\item (\textit{AHSME 1951}) $A$ can do a piece of work in $9$ days. $B$ is $50\%$ more efficient than $A$. The number of days it takes $B$ to do the same piece of work is:


 $\textbf{(A)}\ 13\frac {1}{2} \qquad\textbf{(B)}\ 4\frac {1}{2} \qquad\textbf{(C)}\ 6 \qquad\textbf{(D)}\ 3 \qquad\textbf{(E)}\ \text{none of these answers}$

 \textbf{Answer: }C



\textbf{Solution}

Because $B$ is $50\%$ more efficient, he can do $1.5$ pieces of work in $9$ days. This is equal to 1 piece of work in $\textbf{(C)}\ 6$ days

 


	\item (\textit{AHSME 1951}) In connection with proof in geometry, indicate which one of the following statements is \textit{incorrect}: 


 $\textbf{(A)}\ \text{Some statements are accepted without being proved.}$
$\textbf{(B)}\ \text{In some cases there is more than one correct order in proving certain propositions.}$
$\textbf{(C)}\ \text{Every term used in a proof must have been defined previously.}$
$\textbf{(D)}\ \text{It is not possible to arrive by correct reasoning at a true conclusion if, in the given, there is an untrue proposition.}$
$\textbf{(E)}\ \text{Indirect proof can be used whenever there are two or more contrary propositions.}$

 \textbf{Answer: }E



\textbf{Solution}

$\boxed{\textbf{(C)}}$ is false, as in geometry, the undefined terms point, line, and plane exist.

 A is true, as in geometry, such statements (usually referred to as postulates or axioms) exist.

 B is true, as in geometry, statements may be independent of each other, so it does not matter which order they are proved in. 

 D is true, as if the given or hypothesis of a theorem is false, then that theorem cannot be used.

 E is true, as if two or more propositions are contrary, a proof by contradiction may be used to prove that not all of these propositions can be true.

 


	\item (\textit{AHSME 1951}) The largest number by which the expression $n^3-n$ is divisible for all possible integer values of $n$, is: 


 $\textbf{(A)}\ 2\qquad\textbf{(B)}\ 3\qquad\textbf{(C)}\ 4\qquad\textbf{(D)}\ 5\qquad\textbf{(E)}\ 6$

 \textbf{Answer: }E



\textbf{Solution 1}

Factoring the polynomial gives $(n+1)(n)(n-1)$ According to the factorization, one of those factors must be a multiple of two because there are more than 2 consecutive integers. In addition, because there are three consecutive integers, one of the integers must be a multiple of 3.  Therefore $6$ must divide the given expression.Plugging in $n=2$ yields $6$.  So the largest possibility is $6$.

 Clearly the answer is $\boxed{\textbf{(E)} \ 6}$




\textbf{Solution 2}

In general, $r!$ | $n(n+1)(n+2)...(n+r-1)$ were $r$ and $n$ are integers. So here $3!$ | $n^3$ - $n$ always for any integer $n$.Hence,the correct answer is $6$. $\boxed{\textbf{(E)} \ 6}$



	\item (\textit{AHSME 1951}) If in applying the quadratic formula to a quadratic equation


 
 \[f(x) \equiv ax^2 + bx + c = 0,\]


 it happens that $c = \frac{b^2}{4a}$, then the graph of $y = f(x)$ will certainly:


 $\textbf{(A)}\ \text{have a maximum}\qquad\textbf{(B)}\ \text{have a minimum}\qquad\textbf{(C)}\ \text{be tangent to the x-axis}\\ \qquad\textbf{(D)}\ \text{be tangent to the y-axis}\qquad\textbf{(E)}\ \text{lie in one quadrant only}$

 \textbf{Answer: }C



\textbf{Solution}

The \href{https://artofproblemsolving.com/wiki/index.php?title=Discriminant}{discriminant} of $f(x)$ is $b^2 - 4ac = b^2 - 4a\left(\frac{b^2}{4a}\right) = 0$, so $f(x)$ has exactly one root (with multiplicity 2). At this root, $f(x)$ will therefore $\boxed{\textbf{(C)}\ \text{be tangent to the }x\text{-axis}}$.

 A is false as if $a$ is positive, $f(x)$ will have a minimum, not a maximum.

 B is false as if $a$ is negative, $f(x)$ will have a maximum, not a minimum.

 D is false as $f(x)$ is continuous for all real values of $x$, so it will cross through the $y$-axis instead of being tangent.

 E is false as $f(x)$ is defined for all real values of $x$, so it must lie in at least two quadrants.

 


	\item (\textit{AHSME 1951}) Indicate in which one of the following equations $y$ is neither directly nor inversely proportional to $x$: 


 $\textbf{(A)}\ x+y = 0\qquad\textbf{(B)}\ 3xy = 10\qquad\textbf{(C)}\ x = 5y \qquad\textbf{(D)}\ 3x+y = 10 \qquad\textbf{(E)}\ x/y = \sqrt{3}\qquad$

 \textbf{Answer: }D



\textbf{Solution}

Notice that for any directly or inversely proportional values, it can be expressed as $\frac{x}{y}=k$ or $xy=k$. Now we try to convert each into its standard form counterpart.

 $\textbf{(A)}\ x + y = 0\implies \frac{x}{y}=-1$
$\textbf{(B)}\ 3xy = 10\implies xy=\frac{10}{3}$
$\textbf{(C)}\ x = 5y\implies \frac{x}{y}=5$
$\textbf{(E)}\ \frac {x}{y} = \sqrt {3}\implies \frac {x}{y} = \sqrt {3}$

 As we can see, the only equation without a "standard" form is $\textbf{(D)}$, so our answer is $\boxed{\textbf{(D)}\ 3x + y = 10}$



	\item (\textit{AHSME 1951}) The expression $21x^2+ax+21$ is to be factored into prime binomial factors and without a numerical monomial factor. This can be done if the value ascribed to $a$ is:


 $\textbf{(A)}\ \text{any odd number}\qquad\textbf{(B)}\ \text{some odd number}\qquad\textbf{(C)}\ \text{any even number}\qquad\textbf{(D)}\ \text{some even number}\qquad\textbf{(E)}\ \text{zero}$

 \textbf{Answer: }D



\textbf{Solution 1}

The factors of $21x^2+ax+21$ must be in the format $(bx+c)(dx+e)$, where $b$, $c$, $d$, and $e$ are integers, $bd = ce = 21$ and $cd + be = a$.  Since $b$, $c$, $d$, and $e$ are integer factors of $21$, they must all be odd, and thus $a = cd + be$ must be even.  Also, there are only finitely many possibilities for $(b, c, d, e)$, so there are only finitely many possibilities for $a$. One possible value is $a = 58$, as $21x^2 + 58x + 21 = (7x+3)(3x+7)$. Therefore, $a$ could be $\boxed{\textbf{(D)}\ \text{some even number}}$




\textbf{Solution 2}

Factoring $21x^2+ax+21$ by grouping, we need to find some $b,c$ such that $b\cdot c = 21^2$, and that $b+c=a$.Since $21^2\equiv 1\;(mod\;2)$, $b\land c \equiv 1\;(mod\;2)$, and $b+c \equiv 0\;(mod\;2)$. So $a$ must be even. $a$ cannot be $any$ even number, since $21^2$ only has 4 odd factors, so the answer is $\textbf{(D)}\ \text{some even number}$



	\item (\textit{AHSME 1951}) A six place number is formed by repeating a three place number; for example, $256256$ or $678678$, etc. Any number of this form is always exactly divisible by: 


 $\textbf{(A)}\ 7\text{ only}\qquad\textbf{(B)}\ 11\text{ only}\qquad\textbf{(C)}\ 13\text{ only}\qquad\textbf{(D)}\ 101\qquad\textbf{(E)}\ 1001$

 \textbf{Answer: }E



\textbf{Solution}

We can express any of these types of numbers in the form $\overline{abc}\times 1001$, where $\overline{abc}$ is a 3-digit number. Therefore, the answer is $\textbf{(E)}\ 1001$.

 



\textbf{Solution 2}

$\overline{abcabc}\  = \overline{abc}\* 1001$  .

 


	\item (\textit{AHSME 1951}) When simplified and expressed with negative exponents, the expression $(x+y)^{-1}(x^{-1}+y^{-1})$ is equal to: 


 $\textbf{(A)}\ x^{-2}+2x^{-1}y^{-1}+y^{-2}\qquad\textbf{(B)}\ x^{-2}+2^{-1}x^{-1}y^{-1}+y^{-2}\qquad\textbf{(C)}\ x^{-1}y^{-1}$
$\textbf{(D)}\ \text{some even number}\qquad\textbf{(E)}\ \text{zero}$

 \textbf{Answer: }C



\textbf{Solution}

Note that $(x + y)^{-1}(x^{-1} + y^{-1}) = \dfrac{1}{x + y}\cdot\left(\dfrac{1}{x} + \dfrac{1}{y}\right) = \dfrac{1}{x + y}\cdot\dfrac{x + y}{xy} = \dfrac{1}{xy} = x^{-1}y^{-1}$. The answer is $\textbf{(C)}$.

 


	\item (\textit{AHSME 1951}) Given: $x > 0, y > 0, x > y$ and $z\not = 0$. The inequality which is not always correct is: 


 $\textbf{(A)}\ x+z > y+z\qquad\textbf{(B)}\ x-z > y-z\qquad\textbf{(C)}\ xz > yz$
$\textbf{(D)}\ \frac{x}{z^{2}}>\frac{y}{z^{2}}\qquad\textbf{(E)}\ xz^{2}> yz^{2}$

 \textbf{Answer: }C



\textbf{Solution}

$\textbf{(A)}\ x + z > y + z\implies x>y$, just subtract $z$ from both sides

 $\textbf{(B)}\ x - z > y - z\implies x>y$, just add $z$ to both sides

 $\textbf{(C)}\ xz > yz\implies x>y\text{ if }x>0$, so that means that our desired answer is $\boxed{\textbf{(C)}\ xz > yz}$. 

 As a check:

 $\textbf{(D)}\ \frac {x}{z^2} > \frac {y}{z^2}\implies x>y$, we can divide $z^2$ safely and without worry because $z^2>0$.

 $\textbf{(E)}\ xz^2 > yz^2\implies x>y$, similar reasoning as above but instead, multiply $z^2$.

 


	\item (\textit{AHSME 1951}) The values of $a$ in the equation: $\log_{10}(a^{2}-15a) = 2$ are:


 $\textbf{(A)}\ \frac{15\pm\sqrt{233}}{2}\qquad\textbf{(B)}\ 20,-5\qquad\textbf{(C)}\ \frac{15\pm\sqrt{305}}{2}\qquad\textbf{(D)}\ \pm20$
$\textbf{(E)}\ \text{none of these}$

 \textbf{Answer: }B



\textbf{Solution}

Putting into exponential form, we get that $10^2=a^2-15a\Rightarrow a^2-15a-100=0$

 Now we use the quadratic formula to solve for $a$, and we get $a=\frac{15\pm\sqrt{625}}{2}\implies a=\boxed{\textbf{(B)}\ 20, - 5}$



	\item (\textit{AHSME 1951}) The radius of a cylindrical box is $8$ inches and the height is $3$ inches. The number of inches that may be added to either the radius or the height to give the same nonzero increase in volume is: 


 $\textbf{(A)}\ 1\qquad\textbf{(B)}\ 5\frac{1}{3}\qquad\textbf{(C)}\ \text{any number}\qquad\textbf{(D)}\ \text{non-existent}\qquad\textbf{(E)}\ \text{none of these}$

 \textbf{Answer: }B



\textbf{Solution}

Let $x$ be the number of inches increased. We can set up an equation for $x$: 
 \[8^2 \times (3+x)=(8+x)^2\times 3\]

 Expanding gives $3x^2+48x+192=64x+192$.

 Combining like terms gives the quadratic $3x^2-16x=0$

 Factoring out an $x$ gives $x(3x-16)=0$.

 So either $x=0$, or $3x-16=0$.

 The first equation is not possible, because the problem states that the value has to be non-zero. The second equation gives the answer, $x = \boxed{5\frac{1}{3}}.$ Thus the answer is $\textbf{B}$.

 


	\item (\textit{AHSME 1951}) $\frac{2^{n+4}-2(2^{n})}{2(2^{n+3})}$ when simplified is:


 $\textbf{(A)}\ 2^{n+1}-\frac{1}{8}\qquad\textbf{(B)}\ -2^{n+1}\qquad\textbf{(C)}\ 1-2^{n}\qquad\textbf{(D)}\ \frac{7}{8}\qquad\textbf{(E)}\ \frac{7}{4}$

 \textbf{Answer: }D



\textbf{Solution}

We have $2(2^n)=2^{n+1}$, and $2(2^{n+3})=2^{n+4}$. Thus, $\frac{2^{n+4}-2(2^{n})}{2(2^{n+3})}=\dfrac{2^{n+4}-2^{n+1}}{2^{n+4}}$. Factoring out a $2^{n+1}$ in the numerator, we get $\dfrac{2^{n+1}(2^3-1)}{2^{n+4}}=\dfrac{8-1}{2^3}=\boxed{\textbf{(D)}\ \frac{7}{8}}$.

 


	\item (\textit{AHSME 1951}) The apothem of a square having its area numerically equal to its perimeter is compared with the apothem of an equilateral triangle having its area numerically equal to its perimeter. The first apothem will be: 


 $\textbf{(A)}\ \text{equal to the second}\qquad\textbf{(B)}\ \frac{4}{3}\text{ times the second}\qquad\textbf{(C)}\ \frac{2}{\sqrt{3}}\text{ times the second}\\ \textbf{(D)}\ \frac{\sqrt{2}}{\sqrt{3}}\text{ times the second}\qquad\textbf{(E)}\ \text{indeterminately related to the second}$

 \textbf{Answer: }A



\textbf{Solution 1}

First we try to find the size of the square. Let $s$ be the side length of the square. It states that $s^2=4s$, therefore $s=0,4$. We cross out the trivial case $s=0$, so the side length of the square is $4$. The apothem of the square is simply half its side length, or $2$.

 Let the side length of the equilateral triangle be $t$. The problem states that $3t=\frac{t^2\sqrt{3}}{4}$, so $12t=t^2\sqrt{3}$, therefore $t=0, \frac{12}{\sqrt{3}}$. Again we cross out the trivial case $t=0$, so we have the side length of the triangle as $\frac{12}{\sqrt{3}}$. The apothem of the triangle is $\frac{1}{3}$ of the height. The height of the triangle is $\frac{\sqrt{3}}{2}\cdot \frac{12}{\sqrt{3}} = 6$, so the apothem is $2$ Therefore, the answer is $\textbf{(A)}\ \text{equal to the second}$




\textbf{Solution 2}

Recall that the area of any regular polygon with apothem $a$ and perimeter $p$ is $\frac{1}{2}ap$.  If $p = \frac{1}{2}ap$, then $a = 2$, so any regular polygon whose perimeter is numerically equal to its area has an apothem of length $2$ and $\boxed{\textbf{(A)}}$ is correct.

 


	\item (\textit{AHSME 1951}) In the equation $\frac{x(x-1)-(m+1)}{(x-1)(m-1)}=\frac{x}{m}$ the roots are equal when:


 $\textbf{(A)}\ m = 1\qquad\textbf{(B)}\ m =\frac{1}{2}\qquad\textbf{(C)}\ m = 0\qquad\textbf{(D)}\ m =-1\qquad\textbf{(E)}\ m =-\frac{1}{2}$

 \textbf{Answer: }E



\textbf{Solution}

Multiplying both sides by $(x-1)(m-1)m$ gives us
 \[xm(x-1)-m(m+1)=x(x-1)(m-1)\]
 \[x^2m-xm-m^2-m=x^2m-xm-x^2+x\]
 \[-m^2-m=-x^2+x\]
 \[x^2-x-(m^2+m)=0\]The roots of this quadratic are equal if and only if its discriminant($b^2-4ac$) evaluates to 0.This means
 \[(-1)^2-4(-m^2-m)=0\]
 \[4m^2+4m+1=0\]
 \[(2m+1)^2=0\]
 \[m=-\frac{1}{2} \Rightarrow \boxed{\textbf{(E)}}\]



	\item (\textit{AHSME 1951}) Through a point inside a triangle, three lines are drawn from the vertices to the opposite sides forming six triangular sections. Then: 


 $\textbf{(A)}\ \text{the triangles are similar in opposite pairs}\qquad\textbf{(B)}\ \text{the triangles are congruent in opposite pairs}$
$\textbf{(C)}\ \text{the triangles are equal in area in opposite pairs}\qquad\textbf{(D)}\ \text{three similar quadrilaterals are formed}$
$\textbf{(E)}\ \text{none of the above relations are true}$

 \textbf{Answer: }E



\textbf{Solution}

Say we draw 3 different types of triangles as the following

 1.An equilateral triangle with side lengths $x$

 2.An isosceles triangle with side lengths $x+3,x+3,x-1$

 3.A scalene triangle with sides $x+7,x-3,x+4$

 1. Gives us 6 congruent $30-60-90$ triangles, which means

 $\textbf{(A)},\textbf{(B)},\textbf{(C)},\textbf{(D)}$ are all true.therefore 1. does not give any unique criteria.

 $\textbf{(A)}$-Since the triangles are congruent with each other, they are similar by a ratio of $1$.

 $\textbf{(B)}$-Since all the triangles are congruent they must always be congruent in any way.

 $\textbf{(C)}$-The triangles are all congruent meaning that they always have equal area.

 $\textbf{(D)}$-The top $2$ triangles, bottom-right $2$ triangles, and bottom-left $2$ triangles all make 3 identical congruent quadrilaterals.

 2. Gives us 3 pairs of adjacent congruent triangles, which means

 $\textbf{(E)}$ is the only choice that works.

 $\textbf{(E)}$-All of the relationships don't work so we must choose the last choice, no relations.

 3. Is not needed as we have found that all of the relationships given are not necessary all the time, so $\textbf(E)$ is the only choice that works.

 So the answer is  $\fbox{\textbf{(E)}}$



	\item (\textit{AHSME 1951}) The pressure $(P)$ of wind on a sail varies jointly as the area $(A)$ of the sail and the square of the velocity $(V)$ of the wind. The pressure on a square foot is $1$ pound when the velocity is $16$ miles per hour. The velocity of the wind when the pressure on a square yard is $36$ pounds is: 


 $\textbf{(A)}\ 10\frac{2}{3}\text{ mph}\qquad\textbf{(B)}\ 96\text{ mph}\qquad\textbf{(C)}\ 32\text{ mph}\qquad\textbf{(D)}\ 1\frac{2}{3}\text{ mph}\qquad\textbf{(E)}\ 16\text{ mph}$

 \textbf{Answer: }C



\textbf{Solution}

Because $P$ varies jointly as $A$ and $V^2$, that means that there is a number $k$ such that $P=kAV^2$. You are given that $P=1$ when $A=1$ and $V=16$. That means that $1=k(1)(16^2) \rightarrow k=\frac{1}{256}$. Then, substituting into the original equation with $P=36$ and $A=9$ (because a square yard is $9$ times a square foot), you get $4=\frac{1}{256}(V^2)$. Solving for $V$, we get $V^2=1024$, so $V=32$. Hence, the answer is $\boxed{C}$.

 


	\item (\textit{AHSME 1951}) Of the following sets of data the only one that does not determine the shape of a triangle is: 


 $\textbf{(A)}\ \text{the ratio of two sides and the inc{}luded angle}\\ \qquad\textbf{(B)}\ \text{the ratios of the three altitudes}\\ \qquad\textbf{(C)}\ \text{the ratios of the three medians}\\ \qquad\textbf{(D)}\ \text{the ratio of the altitude to the corresponding base}\\ \qquad\textbf{(E)}\ \text{two angles}$

 \textbf{Answer: }D



\textbf{Solution}

The answer is $\boxed{\textbf{(D)}}$. The ratio of the altitude to the base is insufficient to determine the shape of a triangle; you also need to know the ratio of the two segments into which the altitude divides the base.

 


	\item (\textit{AHSME 1951}) If two poles $20''$ and $80''$ high are $100''$ apart, then the height of the intersection of the lines joining the top of each pole to the foot of the opposite pole is: 


 $\textbf{(A)}\ 50''\qquad\textbf{(B)}\ 40''\qquad\textbf{(C)}\ 16''\qquad\textbf{(D)}\ 60''\qquad\textbf{(E)}\ \text{none of these}$

 \textbf{Answer: }C



\textbf{Solution 1}

The \href{https://artofproblemsolving.com/wiki/index.php?title=Two\_poles\_formula}{two poles formula} says this height is half the harmonic mean of the heights of the two poles. (The distance between the poles is irrelevant.) So the answer is $\frac1{\frac1{20}+\frac1{80}}$, or $\frac1{\frac1{16}}=\boxed{16 \textbf{ (C)}}$.

 



\textbf{Solution 2}

The two lines can be represented as $y=\frac{-x}{5}+20$ and $y=\frac{4x}{5}$.Solving the system,

 $\frac{-x}{5}+20=\frac{4x}{5}$
$20=x.$

 So the lines meet at an $x$-coordinate of 20.

 Solving for the height they meet,

 
 \[y=\frac{4\cdot 20}{5}\]
 \[y=\boxed{16 \textbf{ (C)}}.\]



	\item (\textit{AHSME 1951}) A total of $28$ handshakes were exchanged at the conclusion of a party. Assuming that each participant was equally polite toward all the others, the number of people present was: 


 $\textbf{(A)}\ 14\qquad\textbf{(B)}\ 28\qquad\textbf{(C)}\ 56\qquad\textbf{(D)}\ 8\qquad\textbf{(E)}\ 7$

 \textbf{Answer: }D



\textbf{Solution}

The handshake equation is $h=\frac{n(n-1)}{2}$, where $n$ is the number of people and $h$ is the number of handshakes. There were $28$ handshakes, so $28=\frac{n(n-1)}{2}$$56=n(n-1)$The factors of $56$ are:$1, 2, 4, 7, 8, 14, 28, 56$.As we can see, only $7, 8$ fit the requirements $n(n-1)$ if $n$ was an integer.Therefore, the answer is $\textbf{(D)}\ 8$



	\item (\textit{AHSME 1951}) If $\triangle ABC$ is inscribed in a semicircle whose diameter is $AB$, then $AC+BC$ must be 


 $\textbf{(A)}\ \text{equal to }AB\qquad\textbf{(B)}\ \text{equal to }AB\sqrt{2}\qquad\textbf{(C)}\ \geq AB\sqrt{2}\qquad\textbf{(D)}\ \leq AB\sqrt{2} \qquad \textbf{(E)}\ AB^{2}$

 \textbf{Answer: }D



\textbf{Solution}

Because $AB$ is the diameter of the semi-circle, it follows that $\angle C = 90$. Now we can try to eliminate all the solutions except for one by giving counterexamples.

 $\textbf{(A):}$ Set point $C$ anywhere on the perimeter of the semicircle except on $AB$. By triangle inequality, $AC+BC>AB$, so $\textbf{(A)}$ is wrong.$\textbf{(B):}$Set point $C$ on the perimeter of the semicircle infinitesimally close to $AB$, and so $AC+BC$ almost equals $AB$, therefore $\textbf{(B)}$ is wrong.$\textbf{(C):}$Because we proved that $AC+BC$ can be very close to $AB$ in case $\textbf{(B)}$, it follows that $\textbf{(C)}$ is wrong.$\textbf{(E):}$Because we proved that $AC+BC$ can be very close to $AB$ in case $\textbf{(B)}$, it follows that $\textbf{(E)}$ is wrong.Therefore, the only possible case is $\boxed{\textbf{(D)}\ \leq AB\sqrt{2}}$



	\item (\textit{AHSME 1951}) The roots of the equation $x^{2}-2x = 0$ can be obtained graphically by finding the abscissas of the points of intersection of each of the following pairs of equations except the pair: 


 $\textbf{(A)}\ y = x^{2}, y = 2x\qquad\textbf{(B)}\ y = x^{2}-2x, y = 0\qquad\textbf{(C)}\ y = x, y = x-2\qquad\textbf{(D)}\ y = x^{2}-2x+1, y = 1$
$\textbf{(E)}\ y = x^{2}-1, y = 2x-1$


 \textit{[Note: Abscissa means x-coordinate.]}


 
 \textbf{Answer: }C



\textbf{Solution}

If you find the intersections of the curves listed in the answers $\textbf{(A)}$, $\textbf{(B)}$, $\textbf{(D)}$, and $\textbf{(E)}$, you will find that their abscissas are $0$ and $2$. Also you can note that the curves in $\textbf{(C)}$ don't actually intersect. Therefore the answer is $\boxed{\textbf{(C)}}$.

 


	\item (\textit{AHSME 1951}) The value of $10^{\log_{10}7}$ is: 


 $\textbf{(A)}\ 7\qquad\textbf{(B)}\ 1\qquad\textbf{(C)}\ 10\qquad\textbf{(D)}\ \log_{10}7\qquad\textbf{(E)}\ \log_{7}10$

 \textbf{Answer: }A



\textbf{Solution}

$\log_{10}7=x \Rightarrow 10^x=7$Substitute $\log_{10}7=x$ in $10^{\log_{10}7} \Rightarrow 10^x=?$It was already stated that $10^x=7$, so our answer is $\textbf{(A)}\ 7$



	\item (\textit{AHSME 1951}) If $a^{x}= c^{q}= b$ and $c^{y}= a^{z}= d$, then


 $\textbf{(A)}\ xy = qz\qquad\textbf{(B)}\ \frac{x}{y}=\frac{q}{z}\qquad\textbf{(C)}\ x+y = q+z\qquad\textbf{(D)}\ x-y = q-z$
$\textbf{(E)}\ x^{y}= q^{z}$

 \textbf{Answer: }A



\textbf{Solution}

Try solving both equations for $a$. Taking the $x$-th root of both sides in the first equation and the $z$-the root of both sides in the second gives $a=c^{\frac{q}x}$ and $a=c^{\frac{y}z}$.So $\frac{q}x=\frac{y}z$. Multiplying both sides by $xz$, $qz=xy$. $\boxed{\textbf{(A)}}$.

 


	\item (\textit{AHSME 1951}) Which of the following methods of proving a geometric figure a locus is not correct? 


 $\textbf{(A)}\ \text{Every point of the locus satisfies the conditions and every point not on the locus does not satisfy the conditions.}$
$\textbf{(B)}\ \text{Every point not satisfying the conditions is not on the locus and every point on the locus does satisfy the conditions.}$
$\textbf{(C)}\ \text{Every point satisfying the conditions is on the locus and every point on the locus satisfies the conditions.}$
$\textbf{(D)}\ \text{Every point not on the locus does not satisfy the conditions and every point not satisfying}\\ \text{the conditions is not on the locus.}$
$\textbf{(E)}\ \text{Every point satisfying the conditions is on the locus and every point not satisfying the conditions is not on the locus.}$

 \textbf{Answer: }B



\textbf{Solution}

Statement $\boxed{\textbf{(B)}}$ is wrong because it does not imply that all points that satisfy the conditions are on the locus.

 


	\item (\textit{AHSME 1951}) A number which when divided by $10$ leaves a remainder of $9$, when divided by $9$ leaves a remainder of $8$, by $8$ leaves a remainder of $7$, etc., down to where, when divided by $2$, it leaves a remainder of $1$, is: 


 $\textbf{(A)}\ 59\qquad\textbf{(B)}\ 419\qquad\textbf{(C)}\ 1259\qquad\textbf{(D)}\ 2519\qquad\textbf{(E)}\ \text{none of these answers}$

 \textbf{Answer: }D



\textbf{Solution}

If we add $1$ to the number, it becomes divisible by $10, 9, 8, \cdots, 2, 1$. The LCM of $1$ through $10$ is $2520$, therefore the number we want to find is $2520-1=\boxed{\textbf{(D)}\ 2519}$



	\item (\textit{AHSME 1951}) A rise of $600$ feet is required to get a railroad line over a mountain. The grade can be kept down by lengthening the track and curving it around the mountain peak. The additional length of track required to reduce the grade from $3\%$ to $2\%$ is approximately: 


 $\textbf{(A)}\ 10000\text{ ft.}\qquad\textbf{(B)}\ 20000\text{ ft.}\qquad\textbf{(C)}\ 30000\text{ ft.}\qquad\textbf{(D)}\ 12000\text{ ft.}\qquad\textbf{(E)}\ \text{none of these}$

 \textbf{Answer: }A



\textbf{Solution}

A grade is the rise divided by the horizontal length for a given segment of track. This means we can get the horizontal length of the track by dividing the rise by the grade.

 At a $3\%$ grade, the horizontal track length is $20000$ feet. At a $2\%$ grade, the horizontal track length is $30000$ feet. The difference is $10000$ feet of horizontal track. Compared to $10000$, $600$ feet is insignificant and can be safely covered by the problem's use of the word "approximately." Therefore, the answer is $\boxed{\textbf{(A)}}$.

 


	\item (\textit{AHSME 1951}) A stone is dropped into a well and the report of the stone striking the bottom is heard $7.7$ seconds after it is dropped. Assume that the stone falls $16t^2$ feet in t seconds and that the velocity of sound is $1120$ feet per second. The depth of the well is:


 $\textbf{(A)}\ 784\text{ ft.}\qquad\textbf{(B)}\ 342\text{ ft.}\qquad\textbf{(C)}\ 1568\text{ ft.}\qquad\textbf{(D)}\ 156.8\text{ ft.}\qquad\textbf{(E)}\ \text{none of these}$

 \textbf{Answer: }A



\textbf{Solution}

Let $d$ be the depth of the well in feet, let $t_1$ be the number of seconds the rock took to fall to the bottom of the well, and let $t_2$ be the number of seconds the sound took to travel back up the well. We know $t_1+t_2=7.7$. Now we can solve $t_1$ for $d$:
 \[d=16t_1^2\]
 \[\frac{d}{16}=t_1^2\]
 \[t_1=\frac{\sqrt{d}}4\]And similarly $t_2$:
 \[d=1120t_2\]
 \[t_2=\frac{d}{1120}\]So $\frac{d}{1120}+\frac{\sqrt{d}}4-7.7=0$. If we let $u=\sqrt{d}$, this becomes a quadratic. 
 \[\frac{u^2}{1120}+\frac{u}4-7.7=0\]
 \[u=\frac{-\frac14\pm\sqrt{\left(\frac14\right)^2-4(\frac1{1120})(-7.7)}}{\frac2{1120}}\]
 \[u=560\cdot\left(-\frac14\pm\sqrt{\frac1{16}+\frac{7.7}{280}}\right)\]
 \[u=560\cdot\left(-\frac14\pm\sqrt{\frac{36}{400}}\right)\]
 \[u=560\cdot\left(-\frac14\pm\frac3{10}\right)\]We know $u$ is the positive square root of $d$, so we can replace the $\pm$ with a $+$.
 \[u=560\cdot\left(\frac1{20}\right)\]
 \[u=\frac{560}{20}\]
 \[u=28\]Then $d=28^2=784$, and the answer is $\boxed{\textbf{(A)}}$.

 


	\item (\textit{AHSME 1951}) $\left(\frac{(x+1)^{2}(x^{2}-x+1)^{2}}{(x^{3}+1)^{2}}\right)^{2}\cdot\left(\frac{(x-1)^{2}(x^{2}+x+1)^{2}}{(x^{3}-1)^{2}}\right)^{2}$ equals:


 $\textbf{(A)}\ (x+1)^{4}\qquad\textbf{(B)}\ (x^{3}+1)^{4}\qquad\textbf{(C)}\ 1\qquad\textbf{(D)}\ [(x^{3}+1)(x^{3}-1)]^{2}$
$\textbf{(E)}\ [(x^{3}-1)^{2}]^{2}$

 \textbf{Answer: }C



\textbf{Solution}

First, note that we can pull the exponents out of every factor, since they are all squared. This results in$\left(\frac{(x+1)(x^{2}-x+1)}{x^{3}+1}\right)^{4}\cdot\left(\frac{(x-1)(x^{2}+x+1)}{x^{3}-1}\right)^{4}$Now, multiplying the numerators together gives$\left(\frac{x^3+1}{x^3+1}\right)^{4}\cdot\left(\frac{x^3-1}{x^3-1}\right)^{4}$,which simplifies to $\boxed{1\textbf{ (C)}}$.

 


	\item (\textit{AHSME 1951}) The formula expressing the relationship between $x$ and $y$ in the table is:

 \[\begin{tabular}{|c|c|c|c|c|c|}\hline x & 2 & 3 & 4 & 5 & 6\\ \hline y & 0 & 2 & 6 & 12 & 20\\ \hline\end{tabular}\]
$\textbf{(A)}\ y = 2x-4\qquad\textbf{(B)}\ y = x^{2}-3x+2\qquad\textbf{(C)}\ y = x^{3}-3x^{2}+2x$
$\textbf{(D)}\ y = x^{2}-4x\qquad\textbf{(E)}\ y = x^{2}-4$

 \textbf{Answer: }B



\textbf{Solution}

Just plug the $x,y$ pair $(6,20)$ into each of the 5 answer choices:

 (A): $2(6)-4=8\ne20$

 (B): $6^2-3(6)+2=20$

 (C): $6^3-3(6^2)+2(6)=120\ne20$

 (D): $6^2-4(6)=12\ne20$

 (E): $6^2-4=32\ne20$

 The only one that works is $\boxed{\textbf{(B)}}$.

 


	\item (\textit{AHSME 1951}) If $x =\sqrt{1+\sqrt{1+\sqrt{1+\sqrt{1+\cdots}}}}$, then:


 $\textbf{(A)}\ x = 1\qquad\textbf{(B)}\ 0 < x < 1\qquad\textbf{(C)}\ 1 < x < 2\qquad\textbf{(D)}\ x\text{ is infinite}$
$\textbf{(E)}\ x > 2\text{ but finite}$

 \textbf{Answer: }C



\textbf{Solution}

We note that $x^2=1+\sqrt{1+\sqrt{1+\sqrt{1+\cdots}}}=1+x$. By the quadratic formula, $x=\dfrac{1\pm \sqrt{5}}{2}$. Because there are only positive square roots in $x$, $x$ must be positive, thus, it is $\dfrac{1+\sqrt{5}}{2}\approx 1.618$. Thus, $\boxed{\textbf{(C)}\ 1 < x < 2}$.

 


	\item (\textit{AHSME 1951}) Of the following statements, the only one that is incorrect is:


 $\textbf{(A)}\ \text{An inequality will remain true after each side is increased,}$ $\text{ decreased, multiplied or divided (zero excluded) by the same positive quantity.}$
$\textbf{(B)}\ \text{The arithmetic mean of two unequal positive quantities is greater than their geometric mean.}$
$\textbf{(C)}\ \text{If the sum of two positive quantities is given, their product is largest when they are equal.}$
$\textbf{(D)}\ \text{If }a\text{ and }b\text{ are positive and unequal, }\frac{1}{2}(a^{2}+b^{2})\text{ is greater than }[\frac{1}{2}(a+b)]^{2}.$
$\textbf{(E)}\ \text{If the product of two positive quantities is given, their sum is greatest when they are equal.}$

 \textbf{Answer: }E



\textbf{Solution}

The answer is $\boxed{\textbf{(E)}}$.Quite the opposite of statement (E) is true--the sum $a+b$ is \textit{minimized} when $a=b$, but it approaches $\infty$ when one of $a,b$ gets arbitrarily small.

 


	\item (\textit{AHSME 1951}) If $\frac{xy}{x+y}= a,\frac{xz}{x+z}= b,\frac{yz}{y+z}= c$, where $a, b, c$ are other than zero, then $x$ equals: 


 $\textbf{(A)}\ \frac{abc}{ab+ac+bc}\qquad\textbf{(B)}\ \frac{2abc}{ab+bc+ac}\qquad\textbf{(C)}\ \frac{2abc}{ab+ac-bc}$
$\textbf{(D)}\ \frac{2abc}{ab+bc-ac}\qquad\textbf{(E)}\ \frac{2abc}{ac+bc-ab}$

 \textbf{Answer: }E



\textbf{Solution}

Note that $\frac{1}{a}=\frac{1}{x}+\frac{1}{y}$, $\frac{1}{b}=\frac{1}{x}+\frac{1}{z}$, and $\frac{1}{c}=\frac{1}{y}+\frac{1}{z}$. Therefore

 
 \[\frac{1}{x}+\frac{1}{y}+\frac{1}{z}=\frac{\frac{1}{a}+\frac{1}{b}+\frac{1}{c}}{2}\]

 Therefore

 
 \[\frac{1}{x}=\frac{1}{2a}+\frac{1}{2b}+\frac{1}{2c}-\frac{1}{c}=\frac{1}{2a}+\frac{1}{2b}-\frac{1}{2c}\]

 A little algebraic manipulation yields that

 
 \[x=\boxed{\textbf{(E)}\ \frac{2abc}{ac+bc-ab}}\]



	\item (\textit{AHSME 1951}) If you are given $\log 8\approx .9031$ and $\log 9\approx .9542$, then the only logarithm that cannot be found without the use of tables is:


 $\textbf{(A)}\ \log 17\qquad\textbf{(B)}\ \log\frac{5}{4}\qquad\textbf{(C)}\ \log 15\qquad\textbf{(D)}\ \log 600\qquad\textbf{(E)}\ \log .4$

 \textbf{Answer: }A



\textbf{Solution}

While $\log 17 = \log(8 + 9)$, we cannot easily deal with the logarithm of a sum. Furthermore, $17$ is prime, so none of the logarithm rules involving products or differences works. It therefore cannot be found without the use of a table (note: in 1951, calculators were very rare). The correct answer is therefore $\boxed{\textbf{(A)}\ \log 17}$.

 
As for the rest of the cases: 
 \[\log\frac{5}{4} = \log\frac{10}{8} = \log 10 - \log 8 = 1 - \log 8\] can be found; 
 \[\log 15 = \log 3 + \log 5 = \frac{1}{2}\log 3^2 + \log\frac{10}{2} = \frac{1}{2} \log 9 + \log 10 - \log 2 = \frac{1}{2} \log 9 + 1 - \frac{1}{3} \log 8\] can be found; 
 \[\log 600 = \log 100 + \log 6 = 2 + \log 2 + \log 3 = 2 + \frac{1}{3} \log 8 + \frac{1}{2} \log 9\] can be found; and 
 \[\log .4 = \log\frac{4}{10} = \log 4 - \log 10 = 2 \log 2 - 1 = \frac{2}{3} \log 8 - 1\] can be found.

 


	\item (\textit{AHSME 1951}) $AB$ is a fixed diameter of a circle whose center is $O$. From $C$, any point on the circle, a chord $CD$ is drawn perpendicular to $AB$. Then, as $C$ moves over a semicircle, the bisector of angle $OCD$ cuts the circle in a point that always: 


 $\textbf{(A)}\ \text{bisects the arc }AB\qquad\textbf{(B)}\ \text{trisects the arc }AB\qquad\textbf{(C)}\ \text{varies}$
$\textbf{(D)}\ \text{is as far from }AB\text{ as from }D\qquad\textbf{(E)}\ \text{is equidistant from }B\text{ and }C$

 \textbf{Answer: }A



\textbf{Solution}


                    \begin{figure}[H]
                        \centering
                        \includegraphics[width=0.5\textwidth]{Solution/AHSME_I_1951_46_Solution_Figure_1.png}
                    \end{figure}
                    Draw in diameter $CE$. Note that, as it is inscribed in a semicircle, $\bigtriangleup CDE$ always has a right angle at $D$. Since it is given that $AB \perp CD$, $AB \parallel DE$ by congruent corresponding angles. $CP$ bisects $DE$, as well as the arc subtended by $DE$, $\widehat{DPE}$. Because $AB \parallel DE$, $CP$ always $\boxed{\textbf{(A)}\ \text{bisects the arc }AB}$.

 


	\item (\textit{AHSME 1951}) If $r$ and $s$ are the roots of the equation $ax^2+bx+c=0$, the value of $\frac{1}{r^{2}}+\frac{1}{s^{2}}$ is:


 $\textbf{(A)}\ b^{2}-4ac\qquad\textbf{(B)}\ \frac{b^{2}-4ac}{2a}\qquad\textbf{(C)}\ \frac{b^{2}-4ac}{c^{2}}\qquad\textbf{(D)}\ \frac{b^{2}-2ac}{c^{2}}$
$\textbf{(E)}\ \text{none of these}$

 \textbf{Answer: }D



\textbf{Solution 1}

Note that $\frac{1}{r^2}+\frac{1}{s^2} = \frac{r^2+s^2}{r^2s^2} = \frac{(r+s)^2-2(rs)}{(rs)^2}$. 

 By \href{https://artofproblemsolving.com/wiki/index.php?title=Vieta\%27s}{Vieta's}, this is $\frac{(-\frac{b}{a})^2-2(\frac{c}{a})}{(\frac{c}{a})^2} = \frac{\frac{b^2}{a^2}-\frac{2c}{a}}{\frac{c^2}{a^2}} = \frac{b^2-2ac}{c^2} \implies \boxed{(D)}$




\textbf{Solution 2}

$r$ and $s$ can be found in terms of $a$, $b$, and $c$ by using the quadratic formula; the roots are

 
 \[\frac{-b\pm\sqrt{b^2-4ac}}{2a}\]

 By Vieta's Formula, $r+s=-\frac{b}{a}$ and $rs=\frac{c}{a}$. Now let's algebraically manipulate what we want to find:

 
 \[\frac{1}{r^2}+\frac{1}{s^2}=\frac{r^2+s^2}{r^2s^2}=\frac{(r+s)^2-2rs}{(rs)^2}\]

 Plugging in the values for $r+s$ and $rs$ gives

 
 \[\frac{1}{r^2}+\frac{1}{s^2}=\frac{(-b/a)^2-2(c/a)}{(c/a)^2}=\boxed{\frac{b^2-2ca}{c^2}\textbf{(D)}}\]




\textbf{Solution 3}

\begin{verbatim}
here $r$ and $s$ are the roots of the equation $ax^2+bx+c=0$,
\end{verbatim}

\begin{verbatim}
now we can convert this equation with roots $r^2$ and $s^2$.
\end{verbatim}

\begin{verbatim}
let, $y=x^2$ then above equation becomes
\end{verbatim}

\begin{verbatim}

 \[ay+c=-b\sqrt y\] square on both sides we get 
 \[a^2y^2 +(2ac-b^2)y +c^2 =0\]
\end{verbatim}

\begin{verbatim}
Again we can change this equation with roots $\frac{1}{r^2}$ and $\frac{1}{s^2}$ .
\end{verbatim}

let $z=\frac{1}{y}$ then, $a^2\frac{1}{z^2}+\frac{2ac-b^2}{z}+c^2=0$ then

 \begin{verbatim}
	
 \[c^2z^2+z(2ac-b^2)+a^2=0\]
\end{verbatim}

then the sum of roots of the above equation is $\frac{1}{r^2}+\frac{1}{s^2}=\frac{b^2-2ac}{c^2}$
\begin{verbatim}
hence, 
 \[\frac{1}{r^2}+\frac{1}{s^2}=\boxed{\frac{b^2-2ca}{c^2}\textbf{(D)}}\]
\end{verbatim}




	\item (\textit{AHSME 1951}) The area of a square inscribed in a semicircle is to the area of the square inscribed in the entire circle as: 


 $\textbf{(A)}\ 1: 2\qquad\textbf{(B)}\ 2: 3\qquad\textbf{(C)}\ 2: 5\qquad\textbf{(D)}\ 3: 4\qquad\textbf{(E)}\ 3: 5$

 \textbf{Answer: }C



\textbf{Solution}

Let the radius of the circle be $r$. Let $s_1$ be the side length of the square inscribed in the semicircle and $s_2$ be the side length of the square inscribed in the entire circle. For the square in the semicircle, we have $s_1^2 + (\frac{s_1}{2})^2 = r^2 \Rightarrow s_1^2 = \frac{4}{5}r^2$. For the square in the circle, we have $s_2 \sqrt{2} = 2r \Rightarrow s_2^2 = 2r^2$.

 Therefore, $s_1^2 : s_2^2 = \frac{4}{5}r^2 : 2r^2 = \boxed{\textbf{(C)}\ 2: 5}$



	\item (\textit{AHSME 1951}) The medians of a right triangle which are drawn from the vertices of the acute angles are $5$ and $\sqrt{40}$. The value of the hypotenuse is:


 $\textbf{(A)}\ 10\qquad\textbf{(B)}\ 2\sqrt{40}\qquad\textbf{(C)}\ \sqrt{13}\qquad\textbf{(D)}\ 2\sqrt{13}\qquad\textbf{(E)}\ \text{none of these}$

 \textbf{Answer: }D



\textbf{Solution}

We will proceed by coordinate bashing.

 Call the first leg $2a$, and the second leg $2b$ (We are using the double of a variable to avoid any fractions)

 Notice that we want to find $\sqrt{(2a)^2+(2b)^2}$

 Two equations can be written for the two medians: $a^2 + 4b^2 = 40$ and $4a^2+b^2= 25$. 

 Add them together and we get $5a^2+5b^2=65$,

 Dividing by 5 gives $x^2+y^2=13$ 

 Multiplying this by 4 gives $4x^2+4y^2=52\implies (2x)^2+(2y)^2=52$, just what we need to find the hypotenuse. Recall that he hypotenuse is $\sqrt{(2a)^2+(2b)^2}$. The value inside the radical is equal to $52$, so the hypotenuse is equal to $\sqrt{52}=\boxed{\textbf{(D)}\ 2\sqrt{13}}$



	\item (\textit{AHSME 1951}) Tom, Dick and Harry started out on a $100$-mile journey. Tom and Harry went by automobile at the rate of $25$ mph, while Dick walked at the rate of $5$ mph. After a certain distance, Harry got off and walked on at $5$ mph, while Tom went back for Dick and got him to the destination at the same time that Harry arrived. The number of hours required for the trip was: 


 $\textbf{(A)}\ 5\qquad\textbf{(B)}\ 6\qquad\textbf{(C)}\ 7\qquad\textbf{(D)}\ 8\qquad\textbf{(E)}\ \text{none of these answers}$

 \textbf{Answer: }D



\textbf{Solution 1}

Let $d_1$ be the distance (in miles) that Harry traveled on car, and let $d_2$ be the distance (in miles) that Tom backtracked to get Dick. Let $T$ be the time (in hours) that it took the three to complete the journey. We now examine Harry's journey, Tom's journey, and Dick's journey. These yield, respectively, the equations

 
 \[\frac{d_1}{25}+\frac{100-d_1}{5}=T,\]

 \[\frac{d_1}{25}+\frac{d_2}{25}+\frac{100-(d_1-d_2)}{25}=T,\]

 \[\frac{d_1-d_2}{5}+\frac{100-(d_1-d_2)}{25}=T.\]

 We combine these three equations:

 
 \[\frac{d_1}{25}+\frac{100-d_1}{5}=\frac{d_1}{25}+\frac{d_2}{25}+\frac{100-(d_1-d_2)}{25}=\frac{d_1-d_2}{5}+\frac{100-(d_1-d_2)}{25}\]

 After multiplying everything by 25 and simplifying, we get

 
 \[500-4d_1=100+2d_2=100+4d_1-4d_2\]

 We have that $100+2d_2=100+4d_1-4d_2$, so $4d_1=6d_2\Rightarrow d_1=\frac{3}{2}d_2$. This then shows that $500-4d_1=100+\frac{4}{3}d_1\Rightarrow 400=\frac{16}{3}d_1\Rightarrow d_1=75$, which in turn gives that $d_2=50$. Now we only need to solve for $T$:

 
 \[T=\frac{d_1}{25}+\frac{100-d_1}{5}=\frac{75}{25}+\frac{100-75}{5}=3+5=8\]

 The journey took 8 hours, so the correct answer is $\boxed{\textbf{(D)}}$.

 



\textbf{Solution 2}

As before let $d_1$ be the distance that Harry traveled on car, and let $d_2$ be the distance in miles that Tom backtracked to get Dick and let $T$ be the time (in hours) to complete the journey.  

 For Harry we have
 \[T=\frac{d_1}{25} + \frac{100-d_1}{5}\]
 \[T=\frac{500-4d_1}{25}\]

 for Tom we have
 \[T=\frac{d_1}{25}+\frac{d_2}{25}+ \frac{100-d_1+d_2}{25}\]
 \[T=\frac{100+2d_2}{25}\]for Dick we have
 \[T=\frac{d_1-d_2}{5}+\frac{100-(d_1-d_2)}{25}\]
 \[T=\frac{100+4(d_1-d_2)}{25}\]

 our unknowns then are $T,d_1,d_2$We first get rid of $d_1$ by adding the equations for Harry and Dick and divide by 2 to get 
 \[2T=\frac{500-4d_1+100+4d_1-4d_2}{25}\]
 \[T=\frac{300-2d_2}{25}\]now we can add above to Tom and divide by 2 to get rid of $d_2$ and simultaneously solving for $T$
 \[2T=\frac{300-2d_2+100+2d_2}{25}\]
 \[T=\frac{200}{25}=8\]which is $\boxed{\textbf{(D)}}$.

 


\end{enumerate}